\section{Related Work}

\textbf{Training-free test-time adaptation for CLIP.} TDA~\cite{karmanov2024tda} adapts CLIP
at test time with two small feature caches -- a positive cache keyed by the model's own
predicted class, and a negative cache of low-confidence samples used as a correction term --
built and consumed purely from forward passes, with no backpropagation and no labels. Each new
test image is scored by CLIP and optionally admitted to the positive cache by an entropy
threshold; the final logits are a sum of CLIP's own logits and an $\alpha/\beta$-weighted
cache-logit term computed from the currently cached features. This project's
\texttt{CalTDA}, in \texttt{online\_method/cal\_tda.py}, subclasses TDA directly and reuses its cache-update and
cache-logit machinery unmodified; the three contributions in Section~\ref{sec:method}
change only the admission rule, the CLIP/cache combination rule, and a final temperature.

\textbf{The calibration cost of TTA.} \citet{sheng2025illusion} benchmark a broad set of TTA
methods for vision-language models, including TDA, across the same class of few-shot
recognition datasets used here, and report accuracy alongside expected calibration error
(ECE). Their finding motivates this entire project: TTA's forward-pass-only adaptation loop
tends to raise accuracy while \emph{also} raising ECE, because methods that reward and reuse
confident predictions have no mechanism to distinguish confident-and-correct from
confident-and-wrong. Their reported average ECE for TDA rises 5.70\% $\rightarrow$ 9.21\% on
their suite even as accuracy improves. This project targets that specific gap using TDA as the
adapted method and their framing of the problem, while running fully independent
experiments (Section~\ref{sec:experiments}) rather than reusing their reported numbers.

\textbf{Temperature scaling and ECE.} Expected calibration error (ECE) buckets predictions by
confidence and measures the gap between confidence and empirical accuracy within each bucket,
weighted by bucket occupancy~\cite{guo2017calibration}. Temperature scaling divides logits by
a single learned scalar $T$ before the softmax. The operation is rank-preserving (it cannot
change which class a model predicts) but reshapes confidence, and $T$ is normally fit by
minimizing negative log-likelihood on a labelled validation set. This project cannot use a
labelled validation set -- the spec's no-ground-truth constraint applies throughout adaptation -- so
Section~\ref{sec:method} replaces the usual supervised fit with a leave-one-out, label-free
accuracy estimate computed from the method's own cache.
